% GENERATED FILE. Edit canonical lessons, then run npm run generate:books.
% canonical-source-hash: fnv1a64:c391d6464cfc470f
% canonical-lessons: PT-C268-poderia, PT-C268-seria, PT-C268-revisao-condicional, PT-W268-seis-linhas

\chapter{Would, Could, Should --- The Condicional, and Six Lines}
\label{ch:pt-the-condicional}
\hlchaptermodality{268}

\begin{chapteropening}
\textbf{By the end of this chapter:} \emph{I can say what would happen with the condicional, and write six connected lines about a journey that end looking forward.}

Complete PT-W268-seis-linhas: six lines to one reader, the scene in the imperfeito, events in order, a reason, and a last line with espero, seria or se.
\end{chapteropening}

\section[poderia, gostaria, deveria]{poderia, gostaria, deveria — the would of Portuguese}

\label{lesson:PT-C268-poderia}



\begin{quote}
Say \textquotedblleft{}I would like\textquotedblright{}. (\emph{Gostaria.}) Say \textquotedblleft{}I ought to rest\textquotedblright{}, the careful way. (\emph{Deveria descansar.})

Two forms of one tense. This lesson gives it its name and its recipe.
\end{quote}

\begin{grammarlens}[title={Grammar Lens: the whole infinitive, plus -ia}]
This is the \textbf{condicional}, Portuguese's \textquotedblleft{}would\textquotedblright{}. Keep the \textbf{whole} infinitive and add \textbf{-ia} for \textquotedblleft{}I\textquotedblright{}:

\begin{quote}
falar, \textbf{falaria} — I would speak comer, \textbf{comeria} — I would eat partir, \textbf{partiria} — I would leave
\end{quote}

The other persons change only the ending: \textbf{falaria}, \textbf{falarias}, \textbf{falaria}, \textbf{falaríamos}, \textbf{falariam}. Those are the imperfeito's \emph{-ia} endings, but glued to the whole infinitive: \emph{comia}, I was eating; \emph{comeria}, I would eat.

Three common verbs squeeze their infinitive first: \emph{fazer} gives \textbf{faria}, \emph{dizer} \textbf{diria}, \emph{trazer} \textbf{traria}. \emph{Poder}, \emph{gostar} and \emph{dever} follow the plain recipe: \textbf{poderia}, \textbf{gostaria}, \textbf{deveria}.

Three jobs, all of them about softening. A \textbf{polite request}: \emph{Poderia ajudar-me?} — Could you help me? A \textbf{wish}: \emph{Gostaria de viajar.} And \textbf{advice}: \emph{No teu lugar, estudaria.} — In your place, I'd study.

In everyday European speech the imperfeito often stands in for it, as \emph{queria} did at the ticket window: \emph{podia ajudar-me?}, \emph{gostava de viajar}. Both are right. Notice also the \emph{me} after the verb in \emph{ajudar-me}; Brazil puts it first.
\end{grammarlens}

\begin{cousinweb}
Latin had no would-tense. Spoken Latin made one by gluing a past of \emph{habēre}, \textquotedblleft{}to have\textquotedblright{}, onto the infinitive: roughly \textquotedblleft{}I had to speak\textquotedblright{}, which slid into \textquotedblleft{}I would speak\textquotedblright{}. Portuguese and Spanish chose the \textbf{imperfect} of that verb, \emph{habēbat}, the same past that became \textbf{havia}; worn down on the end of a verb, it is the \textbf{-ia} of \emph{falaria}. Italian glued on the perfect instead, which is why its endings do not match. Neither EtymDB nor Lewis \& Short covers a change this late; this is the usual explanation in histories of Romance, given as that. The name \textbf{condicional} is a learned borrowing from Medieval Latin \emph{conditiōnālis} (EtymDB 2.1).
\end{cousinweb}

\subsection*{Your turn}
Say these aloud:

\begin{itemize}
\raggedright
  \item \emph{falaria, falarias, falaria}
  \item the three squeezed ones — \emph{faria, diria, traria}
  \item a polite request — \emph{Poderia ajudar-me?}
  \item advice — \emph{No teu lugar, estudaria.}
  \item an answer that depends — \emph{Depende. Gostaria de ir, mas não posso.}
\end{itemize}

\begin{tcolorbox}[breakable,colback=teal!4,colframe=teal!35!black,title={Before you move on}]
Build \textquotedblleft{}I would eat\textquotedblright{}. (\emph{\textbf{Comeria}} — the whole infinitive, plus \emph{-ia}.) What are the three jobs of the condicional? (A \textbf{polite request}, a \textbf{wish}, \textbf{advice}.) Which Latin past is hidden in the \emph{-ia}? (\emph{\textbf{Habēbat}}, \textquotedblleft{}had\textquotedblright{}, the same one inside \emph{\textbf{havia}}.)
\end{tcolorbox}
\section[seria]{seria — it would be}

\label{lesson:PT-C268-seria}



\begin{quote}
Build \textquotedblleft{}I would speak\textquotedblright{}. (\emph{Falaria.}) Ask politely for help with \emph{poder}. (\emph{Poderia ajudar-me?})
\end{quote}

\subsection*{You'll want to know: seria}
\textbf{seria} — \textquotedblleft{}it would be\textquotedblright{}, \textquotedblleft{}he or she would be\textquotedblright{}.

\textbf{ser} refused the imperfeito's recipe and gave you \emph{era}. Its condicional does not refuse: the whole infinitive plus \textbf{-ia}, \textbf{seria}, \textbf{serias}, \textbf{seria}, \textbf{seríamos}, \textbf{seriam}. \emph{Seria} is the one you will use most, because it judges a situation rather than a person:

\begin{quote}
\emph{Seria difícil.} — It would be difficult. \emph{Seria mais rápido de comboio.} — It would be quicker by train.
\end{quote}

\emph{De comboio}, \textquotedblleft{}by train\textquotedblright{}, and \emph{de autocarro}, \textquotedblleft{}by bus\textquotedblright{}, are how Portugal says how you travel.

\begin{grammarlens}[title={Grammar Lens: seria bom}]
\textbf{Seria bom} + infinitive — \textquotedblleft{}it would be good to, it would be nice to\textquotedblright{}:

\begin{quote}
\emph{Seria bom viajar juntos.} — It would be nice to travel together.
\end{quote}

That is a wish with the person left out, softer than \emph{gostaria} and less private than \emph{espero}. And it sits naturally after a \emph{depende}, as the verdict on each branch:

\begin{quote}
\emph{Depende. De comboio seria mais rápido, mas de autocarro seria mais barato.}
\end{quote}
\end{grammarlens}

\begin{cousinweb}
\textbf{ser} is inherited from Latin \emph{sedēre}, \textquotedblleft{}to sit\textquotedblright{} (EtymDB 2.1), as chapter sixteen showed, and \textbf{-ia} is the worn Latin \emph{habēbat}, \textquotedblleft{}he had\textquotedblright{}, of the last lesson. So \emph{seria} began as something like \textquotedblleft{}had to be\textquotedblright{}. Neither source records forms this late, so that step is the usual explanation rather than a cited one.
\end{cousinweb}

\subsection*{Your turn}
Say these aloud:

\begin{itemize}
\raggedright
  \item \emph{seria, serias, seria}
  \item \emph{Seria mais rápido de comboio.}
  \item \emph{Seria bom viajar juntos.}
  \item a verdict with \emph{o mais} — \emph{Seria a viagem mais bonita do ano.}
  \item two branches, two verdicts — \emph{De comboio seria mais rápido, mas de autocarro seria mais barato.}
\end{itemize}

\begin{tcolorbox}[breakable,colback=teal!4,colframe=teal!35!black,title={Before you move on}]
Does \emph{ser} follow the recipe in the condicional? (\textbf{Yes}: \emph{\textbf{seria}}.) Say \textquotedblleft{}it would be nice to travel together\textquotedblright{}. (\emph{\textbf{Seria bom viajar juntos.}}) How does Portugal say \textquotedblleft{}by train\textquotedblright{}? (\emph{\textbf{De comboio.}})
\end{tcolorbox}
\section[(review)]{Review — if this, then that would}

\label{lesson:PT-C268-revisao-condicional}



\begin{quote}
Say \textquotedblleft{}it would be difficult\textquotedblright{}. (\emph{Seria difícil.}) Then \textquotedblleft{}I could\textquotedblright{}. (\emph{Poderia.})
\end{quote}

\begin{grammarlens}[title={What you've built}]
Two chapters, and a whole way of talking about what has not happened yet. A real branch: \emph{se chove, não saio}. A branch not yet named: \emph{depende}, or \emph{depende de ti}. And the softened result, in the condicional: \emph{poderia}, \emph{deveria}, \emph{seria}.

\begin{quote}
— \emph{Vamos ao parque amanhã?} — \emph{Depende. Seria bom, mas se chover, poderíamos ir ao museu.}
\end{quote}

\emph{Vamos} is \textquotedblleft{}shall we go\textquotedblright{}, from \emph{ir}, and \emph{poderíamos} is \textquotedblleft{}we could\textquotedblright{}, the \emph{nós} form of \emph{poderia}. And listen to \emph{se chover}: one occasion still to come, so the special form you met last chapter.
\end{grammarlens}

\subsection*{Your turn}
Say these aloud:

\begin{itemize}
\raggedright
  \item \emph{Seria bom ir ao parque amanhã.}
  \item \emph{Se chover, poderíamos ir ao museu.}
  \item advice inside a branch — \emph{Se estás constipado, deverias descansar.}
  \item a hope and a verdict — \emph{Espero viajar. Seria bom.}
  \item your own plan for tomorrow, with one \emph{se} and one \emph{seria}
\end{itemize}

\begin{tcolorbox}[breakable,colback=teal!4,colframe=teal!35!black,title={Before you move on}]
In \emph{se chover, poderíamos ir ao museu}, which half is the condicional? (\emph{\textbf{Poderíamos ir.}}) Which form does the \emph{se} half take, and why? (\emph{\textbf{Chover}} — one occasion still to come.) What would \textquotedblleft{}we would be\textquotedblright{} be, from \emph{seria}? (\emph{\textbf{Seríamos.}}) And which two Latin pieces hide in the \emph{dever} of \emph{deverias}? (\emph{\textbf{Dē}} and \emph{\textbf{habēre}}: to have from someone.)
\end{tcolorbox}
\section[(six lines)]{Six lines about a journey}

\label{lesson:PT-W268-seis-linhas}



\begin{quote}
The writing chapters asked for four lines that held together. This one asks for six, and they tell a story.

Pick one person who will read it: \emph{um amigo}, \emph{uma amiga}, someone in your family.
\end{quote}

\subsection*{What to know first}
A told story is not a list of events. It has a \textbf{scene}, in the imperfeito: \emph{era}, \emph{havia}, \emph{fazia frio}, \emph{estava contente}. It has \textbf{events}, in the perfeito, put in order with \emph{primeiro}, \emph{depois} and \emph{no fim}. It has at least one \textbf{reason}, pointing back with \emph{porque} or forward with \emph{por isso}, and perhaps a turn with \emph{mas}.

And a good one ends by looking forward: \emph{espero …}, or \emph{seria bom …}, or a real condition with \emph{se}.

This book's B1 writing paper asks for a connected text of 130 to 170 words. Six lines is the first step towards it, not the whole of it.

\subsection*{Writing — connected composition}
On paper, write \textbf{six lines} to your reader about a journey, real or invented, one line for each of these:

\begin{enumerate}
\raggedright
  \item the scene: when it was, and the weather;
  \item where you were, who was with you, and how long you stayed, with \emph{durante};
  \item the first thing that happened, with \emph{primeiro};
  \item the next thing, with \emph{depois}, and a reason, with \emph{porque} or \emph{por isso};
  \item how it ended and how you felt, with \emph{no fim} and \emph{estava};
  \item one line forward, with \emph{espero}, \emph{seria bom} or \emph{se}.
\end{enumerate}

Use only words you already write. The whole exercise is choosing the right past for each line, and joining them.

\begin{tcolorbox}[breakable,colback=teal!4,colframe=teal!35!black,title={Before you move on}]
Now read your lines as you read the journey in chapter two hundred and sixty. Mark every verb \textbf{C} for \emph{cena}, the scene, or \textbf{E} for \emph{evento}. Lines one and five should be mostly \textbf{C}; three and four mostly \textbf{E}. A line about the weather marked \textbf{E} is the commonest slip: change \emph{choveu} back to \emph{chovia}.

Then two spelling checks. The \emph{nós} form of the perfeito carries an accent in Portugal, \emph{chegámos}, \emph{esperámos}, and so does the imperfeito's, \emph{estávamos}, \emph{tínhamos}. And every til stays on: \emph{não}, \emph{estação}.

Last, read line six alone. If it has a \emph{se}, is the verb after it in the plain present, or in the special form for one occasion still to come? If it has \emph{seria}, is it your verdict on something that has not happened yet? Then it looks forward, and the story is finished.
\end{tcolorbox}
